%============================================================
%  Bd. 23 - Ungerichtete Graphen
%============================================================

\documentclass[main.tex]{subfiles}

\ifSubfilesClassLoaded{
  \BandSetupStandalone{B23}
}{
}

\title{Bd. 23 - Ungerichtete Graphen}
\author{Martin Kunze}
\date{}
\setcounter{file}{23}

\hypersetup{
  pdftitle={Bd. 23 - Ungerichtete Graphen},
  pdfauthor={Martin Kunze}
}

\begin{document}

\title{Bd. 23 - Ungerichtete Graphen}
\author{Martin Kunze}
\date{}
\hypersetup{pageanchor=false}
\maketitle
\hypersetup{pageanchor=true}
\setcounter{tocdepth}{3}
\tableofcontents
\setcounter{file}{23}
\renewcommand{\FormulaBandID}{23}

\chapter{Einleitung}

Ein ungerichteter Graph wird weiterhin durch eine relationale Kantenmenge
\(E\subseteq V\times V\) dargestellt. Das zusätzliche Symmetrieaxiom
identifiziert die beiden Orientierungen einer Kante: Mit \((x,y)\) gehört
auch \((y,x)\) zu \(E\). Diese Darstellung ist mit den gerichteten Graphen
aus Band~22 unmittelbar verträglich und vermeidet eine zweite Theorie
endlicher Wege.

Der vorliegende Band entwickelt darauf aufbauend ungerichtete
Nachbarschaften und induzierte ungerichtete Teilgraphen. Schleifenfreie,
also schlichte ungerichtete Graphen werden in Band~24 behandelt; Band~25
führt den Zusammenhangsbegriff ein. Auf beiden Spezialisierungen baut die
Baumtheorie in Band~26 auf.

\chapter{Ungerichtete Graphen}

\section{Axiomatischer Begriff}

\FormulaDefDeltaK[Begriff des ungerichteten Graphen]{%
  \UndirectedGraphStruct{V,E}%
}{UndirectedGraphDef}{%
  \DeltaRow{Mengen}{V\dsep E\dsep x\dsep y}
  \DeltaRow{\textbf{Axiome}}{}
  \DeltaRow{Grundgraph}{\GraphStruct{V,E}}
    [\FormulaRefAuto{GraphDef}[def]]
  \DeltaRow{Symmetrie}{%
    x,y\in V\dsep(x,y)\in E\vdash(y,x)\in E}
  \DeltaRow{\textbf{Neues Symbol}}{}
  \DeltaPrem{Ungerichteter Graph}{\UndirectedGraphStruct{V,E}}
}

\FormulaAxiomDeltaK[Zugriff auf den gerichteten Grundgraphen]{%
  \UndirectedGraphStruct{V,E}\vdash E\subseteq V\times V%
}{UndirectedGraphGraphAxiom}{%
  \DeltaRow{Mengen}{V\dsep E}%
}

\FormulaAxiomDeltaK[Zugriff auf die Kantensymmetrie]{%
  \UndirectedGraphStruct{V,E}\vdash
  \forall x,y\in V\,\bigl((x,y)\in E\rightarrow(y,x)\in E\bigr)%
}{UndirectedGraphSymmetryAxiom}{%
  \DeltaRow{Mengen}{V\dsep E\dsep x\dsep y}%
}

\FormulaThmDeltaK[Strukturzugriff auf einen ungerichteten Graphen]{%
  \UndirectedGraphStruct{V,E}\vdash\GraphStruct{V,E}%
}{UndirectedGraphAccess}{%
  \DeltaRow{Mengen}{V\dsep E}%
}
\begin{tabproof}
  \proofstep{1}{\UndirectedGraphStruct{V,E}}{\rA}
  \proofstep{1}{E\subseteq V\times V}{%
    \FormulaRefAuto{UndirectedGraphGraphAxiom}[ax]{1}}
  \proofstep{1}{\GraphStruct{V,E}}{%
    \FormulaRefAuto{GraphDef}[def]{2}}
\end{tabproof}

\FormulaThmDeltaK[Einführungskriterium für ungerichtete Graphen]{%
  \begin{aligned}[t]
    &E\subseteq V\times V\dsep
      \forall x,y\in V\,\bigl((x,y)\in E\rightarrow(y,x)\in E\bigr)
      \vdash{}\\[-2pt]
    &\UndirectedGraphStruct{V,E}
  \end{aligned}%
}{UndirectedGraphIntroduction}{%
  \DeltaRow{Mengen}{V\dsep E\dsep x\dsep y}%
}
\begin{tabproofwide}
  \proofstepwidestar[1]{E\subseteq V\times V}{\rA}
  \proofstepwidestar[2]{%
    \forall x,y\in V\,\bigl((x,y)\in E\rightarrow(y,x)\in E\bigr)}{\rA}
  \proofstepwidestar[1]{\GraphStruct{V,E}}{%
    \FormulaRefAuto{GraphDef}[def]{1}}
  \proofstepwidestar[4]{x\in V}{\rA}
  \proofstepwidestar[5]{y\in V}{\rA}
  \proofstepwidestar[6]{(x,y)\in E}{\rA}
  \proofstepwidestar[2,4,5]{(x,y)\in E\rightarrow(y,x)\in E}{%
    \rRE{4,\rRE{5,2}}}
  \proofstepwidestar[2,4,5,6]{(y,x)\in E}{\rRE{6,7}}
  \proofstepwidestar[2,4,5]{(x,y)\in E\rightarrow(y,x)\in E}{%
    \rRI{6,8}}
  \proofstepwidestar[2,4]{%
    \forall y\in V\,\bigl((x,y)\in E\rightarrow(y,x)\in E\bigr)}{%
    \rUI{\rRI{5,9}}}
  \proofstepwidestar[2]{%
    \forall x,y\in V\,\bigl((x,y)\in E\rightarrow(y,x)\in E\bigr)}{%
    \rUI{\rRI{4,10}}}
  \proofstepwidestar[1,2]{\UndirectedGraphStruct{V,E}}{%
    \FormulaRefAuto{UndirectedGraphDef}[def]{3,11}}
\end{tabproofwide}

\FormulaThmDeltaK[Symmetrie der ungerichteten Adjazenz]{%
  \UndirectedGraphStruct{V,E}\dsep x,y\in V\dsep(x,y)\in E
  \vdash(y,x)\in E%
}{UndirectedAdjacencySymmetric}{%
  \DeltaRow{Mengen}{V\dsep E\dsep x\dsep y}%
}
\begin{tabproof}
  \proofstep{1}{\UndirectedGraphStruct{V,E}}{\rA}
  \proofstep{2}{x\in V}{\rA}
  \proofstep{3}{y\in V}{\rA}
  \proofstep{4}{(x,y)\in E}{\rA}
  \proofstep{1}{%
    \forall x,y\in V\,\bigl((x,y)\in E\rightarrow(y,x)\in E\bigr)}{%
    \FormulaRefAuto{UndirectedGraphSymmetryAxiom}[ax]{1}}
  \proofstep{1,2,3,4}{(y,x)\in E}{%
    \rRE{4,\rUE{\rRE{3,\rUE{5}}}}}
\end{tabproof}

\section{Ungerichtete Nachbarschaften}

\FormulaDefDeltaK[Nachbarschaft in einem ungerichteten Graphen]{%
  \UndirectedGraphStruct{V,E}\dsep x\in V\vdash
  \GraphNeighbors{E}{x}\coloneqq\OutNeighbors{E}{x}%
}{UndirectedNeighborhoodDef}{%
  \DeltaRow{Mengen}{V\dsep E\dsep x\dsep\GraphNeighbors{E}{x}}
  \DeltaRow{\textbf{Neues Symbol}}{}
  \DeltaPrem{Nachbarschaft}{\GraphNeighbors{E}{x}}
}

\FormulaThmDeltaK[Elementkriterium der ungerichteten Nachbarschaft]{%
  \UndirectedGraphStruct{V,E}\dsep x,y\in V\vdash
  y\in\GraphNeighbors{E}{x}\leftrightarrow(x,y)\in E%
}{UndirectedNeighborhoodMembership}{%
  \DeltaRow{Mengen}{V\dsep E\dsep x\dsep y}%
}
\begin{tabproofwide}
  \proofstepwidestar[1]{\UndirectedGraphStruct{V,E}}{\rA}
  \proofstepwidestar[2]{x\in V}{\rA}
  \proofstepwidestar[3]{y\in V}{\rA}
  \proofstepwidestar[1]{\GraphStruct{V,E}}{%
    \FormulaRefAuto{UndirectedGraphAccess}[thm]{1}}
  \proofstepwidestar[1,2]{\GraphNeighbors{E}{x}=\OutNeighbors{E}{x}}{%
    \FormulaRefAuto{UndirectedNeighborhoodDef}[def]{1,2}}
  \proofstepwidestar[1,2,3]{%
    y\in\OutNeighbors{E}{x}\leftrightarrow(x,y)\in E}{%
    \FormulaRefAuto{DirectedOutNeighborhoodMembership}[thm]{4,2,3}}
  \proofstepwidestar[1,2,3]{%
    y\in\GraphNeighbors{E}{x}\leftrightarrow(x,y)\in E}{\rLRS{5,6}}
\end{tabproofwide}

\FormulaThmDeltaK[Symmetrie der Nachbarschaft]{%
  \UndirectedGraphStruct{V,E}\dsep x,y\in V\vdash
  y\in\GraphNeighbors{E}{x}\leftrightarrow
  x\in\GraphNeighbors{E}{y}%
}{UndirectedNeighborhoodSymmetric}{%
  \DeltaRow{Mengen}{V\dsep E\dsep x\dsep y}%
}
\begin{tabproofsplitwide}
  \proofpartwideindR[\(\rightarrow\)]{%
    \UndirectedGraphStruct{V,E}\dsep x,y\in V\vdash
    y\in\GraphNeighbors{E}{x}\rightarrow x\in\GraphNeighbors{E}{y}
  }{%
    \UndirectedGraphStruct{V,E}\dsep x,y\in V\vdash
    y\in\GraphNeighbors{E}{x}\rightarrow x\in\GraphNeighbors{E}{y}
  }[UndirectedNeighborhoodSymmetricForward]
    \proofstepwidestar[1]{\UndirectedGraphStruct{V,E}}{\rA}
    \proofstepwidestar[2]{x\in V}{\rA}
    \proofstepwidestar[3]{y\in V}{\rA}
    \proofstepwidestar[4]{y\in\GraphNeighbors{E}{x}}{\rA}
    \proofstepwidestar[1,2,3,4]{(x,y)\in E}{%
      \FormulaRefAuto{P\leftrightarrow Q, P\vdash Q}[thm]{%
        \FormulaRefAuto{UndirectedNeighborhoodMembership}[thm]{1,2,3},4}}
    \proofstepwidestar[1,2,3,4]{(y,x)\in E}{%
      \FormulaRefAuto{UndirectedAdjacencySymmetric}[thm]{1,2,3,5}}
    \proofstepwidestar[1,2,3,4]{x\in\GraphNeighbors{E}{y}}{%
      \FormulaRefAuto{P\leftrightarrow Q, Q\vdash P}[thm]{%
        \FormulaRefAuto{UndirectedNeighborhoodMembership}[thm]{1,3,2},6}}
    \proofstepwidestar[1,2,3]{%
      y\in\GraphNeighbors{E}{x}\rightarrow x\in\GraphNeighbors{E}{y}}{%
      \rRI{4,7}}
  \closeproofpartwideindR

  \proofpartwideindR[\(\leftarrow\)]{%
    \UndirectedGraphStruct{V,E}\dsep x,y\in V\vdash
    x\in\GraphNeighbors{E}{y}\rightarrow y\in\GraphNeighbors{E}{x}
  }{%
    \UndirectedGraphStruct{V,E}\dsep x,y\in V\vdash
    x\in\GraphNeighbors{E}{y}\rightarrow y\in\GraphNeighbors{E}{x}
  }[UndirectedNeighborhoodSymmetricBackward]
    \proofstepwidestar[1]{\UndirectedGraphStruct{V,E}}{\rA}
    \proofstepwidestar[2]{x\in V}{\rA}
    \proofstepwidestar[3]{y\in V}{\rA}
    \proofstepwidestar[1,2,3]{%
      x\in\GraphNeighbors{E}{y}\rightarrow y\in\GraphNeighbors{E}{x}}{%
      \FormulaRefAuto{UndirectedNeighborhoodSymmetricForward}[thm]{1,3,2}}
  \closeproofpartwideindR

  \proofpartwide{Schluss}
    \proofstepwidestar[1]{\UndirectedGraphStruct{V,E}}{\rA}
    \proofstepwidestar[2]{x\in V}{\rA}
    \proofstepwidestar[3]{y\in V}{\rA}
    \proofstepwidestar[1,2,3]{%
      y\in\GraphNeighbors{E}{x}\rightarrow x\in\GraphNeighbors{E}{y}}{%
      \FormulaRefAuto{UndirectedNeighborhoodSymmetricForward}[thm]{1,2,3}}
    \proofstepwidestar[1,2,3]{%
      x\in\GraphNeighbors{E}{y}\rightarrow y\in\GraphNeighbors{E}{x}}{%
      \FormulaRefAuto{UndirectedNeighborhoodSymmetricBackward}[thm]{1,2,3}}
    \proofstepwidestar[1,2,3]{%
      y\in\GraphNeighbors{E}{x}\leftrightarrow
      x\in\GraphNeighbors{E}{y}}{\rLRI{4,5}}
  \closeproofpartwide
\end{tabproofsplitwide}

\chapter{Teilgraphen}

\section{Induzierte ungerichtete Teilgraphen}

\FormulaThmDeltaK[Symmetrie der induzierten Kantenrelation]{%
  \begin{aligned}[t]
    &\UndirectedGraphStruct{V,E}\dsep U\subseteq V\dsep
      x,y\in U\dsep(x,y)\in\DirectedInducedEdges{E}{U}\vdash{}\\[-2pt]
    &(y,x)\in\DirectedInducedEdges{E}{U}
  \end{aligned}%
}{UndirectedInducedAdjacencySymmetric}{%
  \DeltaRow{Mengen}{V\dsep E\dsep U\dsep x\dsep y}%
}
\begin{tabproofwide}
  \proofstepwidestar[1]{\UndirectedGraphStruct{V,E}}{\rA}
  \proofstepwidestar[2]{U\subseteq V}{\rA}
  \proofstepwidestar[3]{x\in U}{\rA}
  \proofstepwidestar[4]{y\in U}{\rA}
  \proofstepwidestar[5]{(x,y)\in\DirectedInducedEdges{E}{U}}{\rA}
  \proofstepwidestar[1]{\GraphStruct{V,E}}{%
    \FormulaRefAuto{UndirectedGraphAccess}[thm]{1}}
  \proofstepwidestar[2,3]{x\in V}{%
    \FormulaRefAuto{A\subseteq B\dsep x\in A\vdash x\in B}[thm]{2,3}}
  \proofstepwidestar[2,4]{y\in V}{%
    \FormulaRefAuto{A\subseteq B\dsep x\in A\vdash x\in B}[thm]{2,4}}
  \proofstepwidestar[1,2,5]{%
    (x,y)\in E\land x\in U\land y\in U}{%
    \FormulaRefAuto{P\leftrightarrow Q, P\vdash Q}[thm]{%
      \FormulaRefAuto{DirectedInducedEdgeMembership}[thm]{6,2},5}}
  \proofstepwidestar[1,2,5]{(x,y)\in E}{\rAEa{9}}
  \proofstepwidestar[1,2,3,4,5]{(y,x)\in E}{%
    \FormulaRefAuto{UndirectedAdjacencySymmetric}[thm]{1,7,8,10}}
  \proofstepwidestar[1,2,3,4,5]{%
    (y,x)\in E\land y\in U\land x\in U}{%
    \FormulaRefAuto{P\dsep Q\dsep R\vdash P\land Q\land R}[thm]{11,4,3}}
  \proofstepwidestar[1,2,3,4,5]{%
    (y,x)\in\DirectedInducedEdges{E}{U}}{%
    \FormulaRefAuto{P\leftrightarrow Q, Q\vdash P}[thm]{%
      \FormulaRefAuto{DirectedInducedEdgeMembership}[thm]{6,2},12}}
\end{tabproofwide}

\FormulaThmDeltaK[Induzierte Teilgraphen bleiben ungerichtet]{%
  \UndirectedGraphStruct{V,E}\dsep U\subseteq V\vdash
  \UndirectedGraphStruct{U,\DirectedInducedEdges{E}{U}}%
}{UndirectedInducedSubgraph}{%
  \DeltaRow{Mengen}{V\dsep E\dsep U}%
}
\begin{tabproofwide}
  \proofstepwidestar[1]{\UndirectedGraphStruct{V,E}}{\rA}
  \proofstepwidestar[2]{U\subseteq V}{\rA}
  \proofstepwidestar[1]{\GraphStruct{V,E}}{%
    \FormulaRefAuto{UndirectedGraphAccess}[thm]{1}}
  \proofstepwidestar[1,2]{\GraphStruct{U,\DirectedInducedEdges{E}{U}}}{%
    \FormulaRefAuto{DirectedInducedSubgraph}[thm]{3,2}}
  \proofstepwidestar[1,2]{%
    \UndirectedGraphStruct{U,\DirectedInducedEdges{E}{U}}}{%
    \FormulaRefAuto{UndirectedGraphDef}[def]{%
      4,\FormulaRefAuto{UndirectedInducedAdjacencySymmetric}[thm]{1,2}}}
\end{tabproofwide}

\end{document}
